%%%%%%%%%%%%%%%%%%%%%%%%%%%%%%%%%%%%%%%%%%%%%%%%%%%%%%%%%%
% FILE: chapter02_arquitectura.tex
%
% PURPOSE
% -------------------------------------------------------
% Explains the internal architecture of the LaTeX template.
%
%
% ARCHITECTURE ROLE
% -------------------------------------------------------
% Technical documentation of the template structure.
%
%
% DEPENDENCIES
% -------------------------------------------------------
% environments.tex
%
%
% NOTES FOR DEVELOPERS
% -------------------------------------------------------
% If the directory structure of the template changes,
% this chapter must be updated accordingly.
%
%%%%%%%%%%%%%%%%%%%%%%%%%%%%%%%%%%%%%%%%%%%%%%%%%%%%%%%%%%


\chapter{Arquitectura del proyecto}

La plantilla está diseñada siguiendo principios de
arquitectura modular inspirados en prácticas de
ingeniería de software.

Cada componente del proyecto tiene una responsabilidad
claramente definida dentro del sistema.



% --------------------------------------------------
% PRINCIPIOS DE ARQUITECTURA
% --------------------------------------------------

\section{Principios de arquitectura}

\begin{definitionbox}[Arquitectura modular]
Una arquitectura modular divide un sistema en componentes
independientes que pueden mantenerse o modificarse de
forma aislada.
\end{definitionbox}

En esta plantilla, el proyecto se divide en varios módulos
que gestionan aspectos distintos del documento.



\begin{itemize}

\item configuración global
\item estilos tipográficos
\item contenido académico
\item recursos gráficos
\item bibliografía
\item acrónimos

\end{itemize}



\begin{notebox}
Separar estas responsabilidades permite modificar el
formato del documento sin alterar el contenido.
\end{notebox}



% --------------------------------------------------
% ESTRUCTURA DEL PROYECTO
% --------------------------------------------------

\section{Estructura del proyecto}

La estructura de carpetas del proyecto es la siguiente.

\begin{lstlisting}[caption={Estructura de directorios del proyecto},label={code:projectstructure}]
latex-university-template/
|-- main.tex
|-- 00_config/
|   |-- config.tex
|   `-- packages.tex
|-- 01_frontmatter/
|-- 02_chapters/
|-- 03_figures/
|-- 04_bibliography/
|-- 05_appendices/
|-- 06_styles/
`-- 07_acronyms/
\end{lstlisting}

\source{Plantilla}

Cada carpeta tiene una función específica dentro del sistema.



% --------------------------------------------------
% ARCHIVO PRINCIPAL
% --------------------------------------------------

\section{Archivo principal}

El archivo principal del proyecto es:

\begin{lstlisting}[caption={Archivo principal del documento},label={code:mainfile}]
main.tex
\end{lstlisting}

\source{Plantilla}

Este archivo controla la estructura general del documento
e incluye los distintos capítulos del trabajo.



% --------------------------------------------------
% FLUJO DE COMPILACIÓN
% --------------------------------------------------

\section{Flujo de compilación}

Durante la compilación del documento, \LaTeX{} sigue
una secuencia similar a la siguiente.

\begin{lstlisting}[caption={Flujo de compilación del documento},label={code:compileflow}]
main.tex

-> carga configuracion global
-> carga paquetes
-> carga estilos tipograficos
-> carga acronimos
-> genera frontmatter
-> compila capitulos
-> genera bibliografia
-> compila apendices
\end{lstlisting}

\source{Plantilla}



% --------------------------------------------------
% RESPONSABILIDAD DE LOS MODULOS
% --------------------------------------------------

\section{Responsabilidad de los módulos}

\begin{itemize}

\item \textbf{00\_config}  
Configuración global del documento.

\item \textbf{01\_frontmatter}  
Portada, dedicatoria, resúmenes.

\item \textbf{02\_chapters}  
Contenido principal.

\item \textbf{03\_figures}  
Imágenes del documento.

\item \textbf{04\_bibliography}  
Base de datos de referencias.

\item \textbf{05\_appendices}  
Apéndices del documento.

\item \textbf{06\_styles}  
Definición de estilos tipográficos.

\item \textbf{07\_acronyms}  
Definición de acrónimos.

\end{itemize}




% --------------------------------------------------
% BENEFICIOS DE LA ARQUITECTURA
% --------------------------------------------------

\section{Beneficios de la arquitectura}

Esta arquitectura modular ofrece varias ventajas.

\begin{itemize}

\item facilita el mantenimiento del proyecto
\item permite reutilizar la plantilla en múltiples documentos
\item separa claramente contenido y formato
\item reduce errores de compilación

\end{itemize}



\begin{examplebox}
Un autor puede escribir todos los capítulos del documento
sin necesidad de modificar los archivos de configuración
o los estilos tipográficos.
\end{examplebox}
